% Титульная страница
% ---------- Авторы: поправьте ФИО здесь ----------
\newcommand{\courseauthors}{%
  Корягин Сергей Андреевич\\
  Полякова Полина Алексеевна}
\newcommand{\authorposition}{педагоги дополнительного образования\\ МОБУ «СОШ «ЦО «Кудрово»}
% =================================================
\begin{titlepage}
\centering
\vspace*{0.6cm}
{\sffamily\Huge\bfseries\color{espdark} ESP32-S3\par}
\vspace{3mm}
{\sffamily\LARGE Первые шаги в электронике\\[1mm] и программировании микроконтроллеров\par}
\vspace{7mm}

% ---------- Иллюстрация: плата-персонаж ----------
\begin{tikzpicture}[x=1cm,y=1cm,line join=round,line cap=round]
  % фон-подложка
  \fill[accent!7,rounded corners=8mm] (-7.2,-3.6) rectangle (7.2,4.9);
  % волна: аналоговый -> цифровой
  \draw[accent,line width=1.6pt,domain=-6.6:-2.8,samples=80] plot (\x,{-2.7+0.45*sin(\x*170)});
  \draw[->,gray,line width=1pt] (-2.6,-2.7) -- (-1.9,-2.7);
  \draw[esp,line width=1.6pt] (-1.7,-3.1) -- (-1.7,-2.3) -- (-1.1,-2.3) -- (-1.1,-3.1) -- (-0.5,-3.1) -- (-0.5,-2.3) -- (0.7,-2.3) -- (0.7,-3.1) -- (1.3,-3.1);
  \node[font=\sffamily\tiny,text=accent] at (-4.7,-3.35) {аналоговый};
  \node[font=\sffamily\tiny,text=esp] at (-0.2,-3.35) {цифровой};
  % плата
  \fill[black!15,rounded corners=5mm] (-2.05,-1.95) rectangle (2.25,3.35);
  \fill[black!82,rounded corners=5mm] (-2.2,-1.8) rectangle (2.1,3.5);
  \foreach \y in {-1.3,-0.8,...,3.1}{
    \fill[pinmetal] (-2.0,\y-0.1) rectangle (-1.8,\y+0.1);
    \fill[pinmetal] (1.7,\y-0.1) rectangle (1.9,\y+0.1);
  }
  % модуль-«лицо»
  \fill[gray!45,rounded corners=3mm] (-1.35,-0.6) rectangle (1.25,2.6);
  \fill[gray!25,rounded corners=3mm] (-1.25,-0.5) rectangle (1.15,2.5);
  % антенна
  \fill[gray!55] (-1.35,2.75) rectangle (1.25,3.3);
  \draw[white,line width=1pt] (-1.1,2.85) -- (-1.1,3.2) -- (-0.75,3.2) -- (-0.75,2.85) -- (-0.4,2.85) -- (-0.4,3.2) -- (-0.05,3.2) -- (-0.05,2.85) -- (0.3,2.85) -- (0.3,3.2) -- (0.65,3.2) -- (0.65,2.85) -- (1.0,2.85);
  % глаза
  \foreach \ex in {-0.55,0.45}{
    \fill[white] (\ex,1.45) ellipse (0.33 and 0.42);
    \fill[black!85] (\ex+0.06,1.38) circle (0.2);
    \fill[white] (\ex+0.12,1.46) circle (0.07);
  }
  % щёки и улыбка
  \fill[esp!45,opacity=0.7] (-0.95,0.72) ellipse (0.22 and 0.13);
  \fill[esp!45,opacity=0.7] (0.85,0.72) ellipse (0.22 and 0.13);
  \draw[black!80,line width=1.4pt] (-0.4,0.7) .. controls (-0.2,0.35) and (0.1,0.35) .. (0.3,0.7);
  \node[font=\sffamily\bfseries\scriptsize,text=black!70] at (-0.05,-0.15) {ESP32-S3};
  % USB-разъёмы и кнопки
  \fill[gray!60,rounded corners=0.5mm] (-1.5,-2.15) rectangle (-0.7,-1.6);
  \fill[gray!60,rounded corners=0.5mm] (0.6,-2.15) rectangle (1.4,-1.6);
  \fill[gray!50] (-1.45,-1.35) rectangle (-1.05,-0.95); \fill[black!60] (-1.25,-1.15) circle (0.12);
  \fill[gray!50] (0.95,-1.35) rectangle (1.35,-0.95); \fill[black!60] (1.15,-1.15) circle (0.12);
  % Wi-Fi дуги
  \foreach \r/\o in {0.55/1,0.95/0.75,1.35/0.5}{
    \draw[okgreen,line width=2pt,opacity=\o] (2.9,3.3) ++(135:\r) arc[start angle=135,end angle=45,radius=\r];
  }
  \fill[okgreen] (2.9,3.3) circle (0.1);
  \node[font=\sffamily\bfseries\scriptsize,text=okgreen!70!black] at (2.9,2.85) {Wi-Fi};
  % провода к светодиоду
  \draw[wyellow,line width=2.2pt,rounded corners=3mm] (1.9,1.6) -- (3.6,1.6) -- (3.6,0.9);
  \draw[black!80,line width=2.2pt,rounded corners=3mm] (1.9,0.1) -- (4.3,0.1) -- (4.3,0.9);
  % светодиод с лучами
  \foreach \a in {30,60,90,120,150}{\draw[warn,line width=1.6pt] (3.95,2.05) ++(\a:0.95) -- ++(\a:0.4);}
  \fill[esp!80!black] (3.35,0.85) rectangle (4.55,1.05);
  \fill[esp] (3.45,1.05) -- (3.45,1.8) arc[start angle=180,end angle=0,radius=0.5] -- (4.45,1.05) -- cycle;
  \fill[white,opacity=0.6] (3.72,1.95) circle (0.1);
  % кнопка и потенциометр слева
  \draw[wgreen,line width=2.2pt,rounded corners=3mm] (-1.8,1.1) -- (-3.4,1.1) -- (-3.4,1.6);
  \fill[black!75,rounded corners=1mm] (-4.0,1.6) rectangle (-2.8,2.8);
  \fill[black!40] (-3.4,2.2) circle (0.38);
  \draw[wblue,line width=2.2pt,rounded corners=3mm] (-1.8,-0.4) -- (-4.9,-0.4) -- (-4.9,0.1);
  \fill[blue!55!black,rounded corners=1mm] (-5.6,0.1) rectangle (-4.2,1.5);
  \fill[white] (-4.9,0.8) circle (0.5);
  \fill[gray!30] (-4.9,0.8) circle (0.38);
  \draw[black!70,line width=1.5pt] (-4.9,0.8) -- ++(60:0.35);
  % экран справа снизу
  \draw[wpurple,line width=2.2pt,rounded corners=3mm] (1.9,-0.9) -- (4.4,-0.9) -- (4.4,-1.2);
  \fill[blue!45!black,rounded corners=1mm] (3.6,-2.9) rectangle (6.4,-1.2);
  \fill[black] (3.8,-2.7) rectangle (6.2,-1.4);
  \node[font=\ttfamily\bfseries\scriptsize,text=cyan!60] at (5.0,-1.8) {Привет!};
  \fill[cyan!60] (3.95,-2.45) rectangle (5.4,-2.25);
\end{tikzpicture}

\vspace{7mm}
\begin{minipage}{0.92\textwidth}
\sffamily\small
\begin{multicols}{2}
\begin{enumerate}[label=\arabic*.,itemsep=0pt]
  \item Знакомство с платой
  \item Сигналы: аналоговые и цифровые
  \item Выводы платы: что есть что
  \item Arduino IDE и первая программа
  \item Светодиод: закон Ома и мигание
  \item ШИМ: яркость и частота
  \item Радио: Wi-Fi и Bluetooth
  \item Веб-сервер: HTTP, HTML и CSS
  \item Последовательный порт
  \item Кнопки и прерывания
  \item АЦП: измеряем напряжение
  \item Шина I\textsuperscript{2}C: экран
  \item FreeRTOS: несколько дел сразу
  \item Сон и энергосбережение
\end{enumerate}
\end{multicols}
\end{minipage}

\vfill
{\sffamily\large\bfseries \courseauthors\par}
\vspace{1mm}
{\sffamily\small \authorposition\par}
\vspace{6mm}
{\footnotesize Все примеры --- для Arduino-ESP32 версии 3.x
и отладочной платы ESP32-S3-DevKitC-1 (или совместимой).\par}
\end{titlepage}
